\slajdid{02-s012}
\begin{frame}[label=02-s012]
\gusttekst
Po pravilu u funkcionalnu tabelu se unose sve moguće vrednosti koje mogu da imaju ulazni signali odnosno logičke promenljive. Ponekad je moguće funkcionalnu tabelu napisati i u „skraćenoj“ formi, ali ona opet suštinski pokriva sve moguće situacije.

Neka je isti digitalni sistem sa tri ulaza i izlaza, ali je zahtev drugačiji.

Izlazni signal je na nivou logičke jedinice ako je bar jedan ulaz na nivou logičke jedinice.

\begin{columns}[c,onlytextwidth]
\begin{column}{.23\textwidth}
\begingroup
\setlength{\tabcolsep}{3pt}
\renewcommand{\arraystretch}{1.2}
\par\noindent
\begin{tabularx}{\linewidth}{*{4}{>{\centering\arraybackslash}X}}
\toprule C&B&A&F\\\midrule
0&0&0&0\\
b&b&1&1\\
b&1&b&1\\
1&b&b&1\\\bottomrule
\end{tabularx}\par
\endgroup

\end{column}
\begin{column}{.72\textwidth}
U ovoj skraćenoj formi upotrebljen je simbol \alert{b – bilo šta}, u smislu da ta logička promenjiva može uzeti i vrednost logičke jedinice i vrednost logičke nule, i da to neće uticati na izlaznu logičku promenljivu.

Ponekad ćemo tu situaciju označavati i sa \alert{X} sa istim značenjem.
\end{column}
\end{columns}
\end{frame}
